% =====================================================================
% GreenStep Challenge Admin -- Project Cover Letter
% CISC 480 Senior Capstone, Spring 2026
% =====================================================================
\documentclass[11pt,letterpaper]{article}

\usepackage[margin=1in]{geometry}
\usepackage{hyperref}
\usepackage{xcolor}
\usepackage{parskip}
\usepackage{microtype}

\hypersetup{
  colorlinks=true,
  linkcolor=blue!60!black,
  urlcolor=blue!60!black,
  pdftitle={GreenStep Admin Console -- Project Cover Letter},
  pdfauthor={Team 2: Josh Dunlap, Eli Goldberger, Khue Vo, Rudy Vergara}
}

\pagestyle{empty}
\setlength{\parskip}{0.6em}
\setlength{\parindent}{0pt}

\begin{document}

% ---- Sender block ----
\begin{flushleft}
\textbf{Team 2 -- GreenStep Capstone} \\
Department of Computer and Information Sciences \\
University of St.\ Thomas \\
2115 Summit Avenue \\
St.\ Paul, MN 55105 \\
\href{https://github.com/CISC-480-GreenStep/sp26_team2_greenstepchallenge_admin}{github.com/CISC-480-GreenStep/sp26\_team2\_greenstepchallenge\_admin}
\end{flushleft}

\vspace{1em}
May 21, 2026

\vspace{1em}

% ---- Recipient block ----
\begin{flushleft}
Kristin Mroz-Risse \\
Sustainability Program Coordinator \\
Minnesota Pollution Control Agency \\
520 Lafayette Road North \\
St.\ Paul, MN 55155

\vspace{0.5em}
And to future student teams \\
Department of Computer and Information Sciences \\
University of St.\ Thomas
\end{flushleft}

\vspace{1.5em}

% ---- Salutation ----
Dear Kristin and future student teams,

% ---- Body ----
This letter accompanies the final deliverables for the GreenStep Challenge Admin
Console, the spring 2026 CISC~480 capstone project at the University of St.\ Thomas.
We are writing both to summarize what we built and to make it as easy as possible
for whoever picks the work up next to keep going.

\textbf{Project summary and impact.} MPCA runs annual sustainability challenges
that ask staff to log everyday actions across food, energy, transportation,
water, and consumption categories. Before this project, that program lived in
spreadsheets and email. The GreenStep Admin Console replaces that workflow with
a single web application backed by a managed Postgres database. The application
lets MPCA administrators design challenges, organize participants into flexible
groups (departments, cohorts, building-level competitions), monitor engagement
on a customizable drag-and-drop dashboard, and export participation reports for
leadership. The console runs at \url{https://greenstep.auriga.fyi} and exposes
three role tiers (SuperAdmin, Admin, GeneralUser) enforced at the database
layer through Postgres Row-Level Security so authorization cannot be bypassed
from the browser.

\textbf{Constraints, limitations, and next steps.} The capstone reached a
working v0.9.0 build, but we want to be transparent about what is \emph{not} in
the box. The standalone Reports page and the new Participation Report dashboard
widget still coexist; consolidating them under a single ``Report and Analysis''
surface is GitHub issue~\#40 and is partially complete. Two stub tables,
\texttt{badges} and \texttt{tickets}, were planned but not finished; their UI
shells exist on a closed pull request (\#80) and are tracked in issue~\#90. The
legacy \texttt{presets} / \texttt{preset\_actions} tables remain on the live
database even though no application code reads from them (issue~\#88). A few
schema-versus-UI inconsistencies around challenge color fields are documented
in issue~\#89. Photo moderation, push notifications, and multi-device dashboard
layout sync are noted as future work in the Project Report. The mobile
application built by Team~1 shares the same database but the two surfaces have
not been exercised together at scale.

\textbf{Implementation and continuity guidance.} The repository is structured
to be readable on first contact. The top-level \texttt{README.md} explains the
project and how to run it; the \texttt{docs/} folder holds the architecture
map, the consolidated product documentation (PRD, UX spec, technical design,
source-code docs, QA activities), and this project report. Coding standards
live in \texttt{CODING\_GUIDELINES.md} and are enforced by ESLint and a
300-line-per-file cap. The database schema is consolidated into two migration
files (\texttt{001\_initial\_schema.sql} and \texttt{002\_seed\_data.sql}); we
squashed fourteen historical migrations into this baseline so the next team
has one place to look, not fourteen.

We would suggest three concrete starting moves for the next team:
(1)~run \texttt{supabase db reset} locally against the consolidated baseline
and confirm the schema matches production;
(2)~walk through issues \#40, \#88, \#89, and \#90 in that order to land the
small hand-off items the team already mapped out; and
(3)~book an early meeting with Kristin to confirm whether the standalone
Reports page should be retired before any new features are added.

\textbf{Gratitude.} Kristin, thank you for the time, candor, and flexibility
you gave us across the semester. The willingness to walk us through how the
agency actually runs these challenges, the patience as our design shifted from
a FastAPI backend to Supabase, and the steady focus on the people the program
is for made this a project we are proud to put our names on. The application
exists because you were willing to treat four student developers as real
collaborators, and we hope the next team gets the same opportunity.

\vspace{1em}

% ---- Closing ----
Sincerely,

\vspace{3em}

% ---- Signature block ----
\begin{tabular}{@{}p{0.45\linewidth}p{0.45\linewidth}@{}}
\rule{0.4\linewidth}{0.4pt} & \rule{0.4\linewidth}{0.4pt} \\
Josh Dunlap & Eli Goldberger \\[1.5em]
\rule{0.4\linewidth}{0.4pt} & \rule{0.4\linewidth}{0.4pt} \\
Khue Vo & Rudy Vergara \\
\end{tabular}

\vspace{1em}
Team 2 -- CISC 480 Senior Capstone, Section 01 (Spring 2026) \\
University of St.\ Thomas

\end{document}
